\section*{Chapter \ref{ch:nw:api-engineering-for-crypto-venues} --- API Engineering for Crypto Venues}

\begin{solution}{exo:nw:api-engineering-for-crypto-venues:1}
$6\,000/250 = 24$, and nothing else in that minute.
\end{solution}

\begin{solution}{exo:nw:api-engineering-for-crypto-venues:2}
To limit what one connection's problems cost: \omterm{def:nw:networking-for-trading:hol}{head-of-line blocking} behind a busy symbol, the books lost when it drops, the snapshots needed to
recover, and the load of the 24-hour reset. Five connections of 160 share the risk and the work.
\end{solution}

\begin{solution}{exo:nw:api-engineering-for-crypto-venues:3}
Answer with a pong carrying the ping's payload, as soon as practical (RFC 6455). The venue pings every 20 seconds and disconnects a connection that has
not answered within a minute.
\end{solution}

\begin{solution}{exo:nw:api-engineering-for-crypto-venues:4}
In 10 seconds one address may open $300 \times 10/300 = 10$ connections on the even-spreading reading, so 40 connections need 4 addresses.
\end{solution}

\begin{solution}{exo:nw:api-engineering-for-crypto-venues:5}
Receive minus event is the one-way delay minus the venue's \omterm{def:nw:time-synchronisation:offset}{clock offset}; the delay has a floor that some messages reach in every window, so the
minimum is the floor minus the offset, and its changes are the offset's. It fails when the delay's floor moves (a route change, congestion lasting the
whole window) or when the venue stamps events at different points for different streams.
\end{solution}

\begin{solution}{exo:nw:api-engineering-for-crypto-venues:6}
It mixes two clocks: the difference includes the venue's \omterm{def:nw:time-synchronisation:offset}{clock offset} (milliseconds in the chapter's simulation) and the one-way delay to the firm.
Measure tick-to-order on the firm's own clock from the arrival timestamp, and use the venue's time only after subtracting the estimated skew.
\end{solution}

\begin{solution}{exo:nw:api-engineering-for-crypto-venues:7}
75: $80 \times 75 = 6\,000$, the whole minute's budget; at 76 the last snapshot waits for the next minute.
\end{solution}

\begin{solution}{exo:nw:api-engineering-for-crypto-venues:8}
One connection is one failure domain: a disconnection or the 24-hour reset drops every book, a busy symbol delays all others, and a per-connection
recovery needs a snapshot for every symbol, which the weight budget cannot pay for quickly. More connections, each below the cap, are more reliable.
\end{solution}

\begin{solution}{pb:nw:api-engineering-for-crypto-venues:1}
\textbf{Part I.}
\begin{enumerate}
\item 800 streams, 5 connections of 160, 1 address.
\item One connection.
\item Still one: five reconnections in 10 seconds are well within the per-address limit.
\item Every connection must be replaced at least daily: stagger the resets and reconnect before the venue does, with a new connection subscribed before
  the old one closes.
\end{enumerate}
\textbf{Part II.}
\begin{enumerate}[resume]
\item Five (the lost messages touch five symbols) per symbol; all 80 of the connection per connection.
\item \qty{180.1}{\second}, and 5\,769 stale symbol-seconds.
\item \qty{0.2}{\second} and 11.2 symbol-seconds: 80 snapshots of weight 50 fit in one minute.
\item Send any request that needs weight, orders included, until the recovery's snapshots have used their windows.
\end{enumerate}
\textbf{Part III.}
\begin{enumerate}[resume]
\item Risk 500 (all it asks), market making 3\,300, arbitrage 2\,200.
\item As a reserved share, or through the risk class: recovery must not be starved by trading, nor starve it.
\item Within \qty{15}{\micro\second} of the truth, from \qty{-2.23}{\milli\second} at the start to \qty{-4.20}{\milli\second} at the end.
\item To measure latency to and from the venue on one time base, and to notice when the venue's clock, and so its timestamps, move.
\end{enumerate}
\textbf{Part IV.}
\begin{enumerate}[resume]
\item \emph{Named result}: 800 streams need five connections of 160 and one address under the venue's limits; five lost messages cost 0.5 stale
  symbol-seconds with per-symbol recovery, and 5\,769 symbol-seconds (the last book back after three minutes) with per-connection recovery from full-depth
  snapshots.
\item Per-symbol recovery; then the snapshot depth.
\item The firm could not tell which symbols missed messages and would have to rebuild every book on the connection: smaller shards and shallower snapshots
  would then be the only defences.
\item Evenly: a few connections at a time, each renewed before the venue closes it.
\item A second copy of every event: a missed message on one is usually present on the other, so the book needs no snapshot; the first copy wins.
\item Fixed windows aligned like the venue's, several rules at once (weight, orders per interval, per day), and the pause on 429 and 418 answers.
\item In the crypto rows: shards, addresses, governors per limit, recovery design and snapshot depth, and the clock-skew monitoring.
\item Either a few tenths of a second of one book, or minutes of every book on a connection and of the firm's whole rate budget.
\end{enumerate}
\end{solution}

\begin{solution}{iq:nw:api-engineering-for-crypto-venues:1}
Buffer the deltas, fetch a snapshot, drop the deltas at or before the snapshot's update id, apply the rest in order; if a delta's first update id is beyond
the book's last plus one, discard the book and start again; verify with the venue's checksum where it offers one.
\iqlookfor{Sequence numbers; the gap rule; snapshot and buffer ordering; checksums.}
\end{solution}

\begin{solution}{iq:nw:api-engineering-for-crypto-venues:2}
Below the venue's per-connection cap (a few hundred streams each), a dozen connections spread over two or more addresses, symbols balanced by message
rate, reconnections staggered, per-symbol recovery, and a spare connection ready.
\iqlookfor{Caps and limits; failure domains; staggering; recovery design.}
\end{solution}

\begin{solution}{iq:nw:api-engineering-for-crypto-venues:3}
Route every request through one governor per venue limit and give each strategy a share by priority (water-filling), with a reserved share for risk and
recovery; refuse or queue what exceeds a share rather than letting the venue answer 429.
\iqlookfor{Central governor; allocation; reserves; no reliance on the venue's errors.}
\end{solution}

\begin{solution}{iq:nw:api-engineering-for-crypto-venues:4}
Stamp arrivals with a synchronised clock and track the minimum of arrival minus the venue's event time over windows: its level is the minimum delay minus the
venue's offset, and its changes are the offset's. Cross-check against round trips measured with the venue's own time endpoint.
\iqlookfor{Minimum-delay filter; separating delay from offset; verification.}
\end{solution}

\begin{solution}{iq:nw:api-engineering-for-crypto-venues:5}
By latency, recovery and limits: a persistent WebSocket or FIX session avoids per-request handshakes and gives sequence numbers (FIX) or push
acknowledgements; REST is simple and stateless but slower; check whether the venue counts all paths against one limit and which offers \omterm{def:nw:private-connectivity-to-crypto-venues:endpoint}{private endpoints}.
\iqlookfor{Session against request; recovery semantics; shared limits.}
\end{solution}

\begin{solution}{iq:nw:api-engineering-for-crypto-venues:6}
The 24-hour reset or a scheduled venue event (connections all opened together, reset together); the machine's own schedule (a cron job, garbage
collection, a backup); message rates at that time (an index rebalance, a funding time); and the limits (5 incoming messages a second, reconnection
limits) hit by a burst of re-subscriptions.
\iqlookfor{Periodic causes; staggering; limits; measure before guessing.}
\end{solution}
